\section{Introduction}\label{sec:introduction}

A partial algebra is an algebra in which some applications of the basic
operations are undefined.  A classical way to compare a partial algebra with
total algebras is to embed it into, or map it to, compatible total algebras.
We study a restricted version of this problem in which the original carrier is
kept pointwise and only finitely many new elements are permitted.

Let $D$ be a one-sorted partial algebra with carrier $A$ and a finite binary
signature $\Sigma$.  We use the standard free compatible completion $\Free D$
of $D$: applications already defined in $D$ are reduced to their prescribed
values, while undefined applications remain as finite formal terms.  The free
completion is background; the question considered here is how finite subsets
of $\Free D$ can be distinguished by compatible total completions of $D$.

\begin{definition}
A \emph{finite-complement completion} of $D$ is a total $\Sigma$-algebra $T$
containing $A$ as a set, such that
\begin{enumerate}
  \item every operation value already defined in $D$ has the same value in
        $T$, and
  \item $T\setminus A$ is finite.
\end{enumerate}
The universal property of $\Free D$ gives a canonical homomorphism
$\pi_T:\Free D\to T$ fixing $A$ pointwise.
\end{definition}

The main result says that the family of all such completions discriminates every
finite subset of the free completion.

\begin{theorem}[Finite-subset discrimination; informal form]
\label{thm:intro-discrimination}
For every finite subset $X\subseteq\Free D$ there exists a finite-complement
completion $T$ of $D$ such that $\pi_T$ is injective on $X$.
\end{theorem}

The proof is elementary and constructive.  Take all subterms occurring in the
chosen finite set.  Keep base elements unchanged, assign one new element to
each non-base subterm that must be remembered, and send all remaining cases to
a single sink.  The normal-form property of the free completion guarantees
that this finite operation table extends the partial operations of $D$.

This statement resembles the standard passage from residual separation to
finite-set discrimination.  There is, however, a simple obstruction to the
usual proof by finite products.  If $A$ is infinite, the class of
finite-complement completions need not be closed under products over the fixed
copy of $A$.  For example, each inclusion
\[
 A\hookrightarrow A\sqcup\{\ast\}
\]
has one new element, whereas the diagonal copy of $A$ in
$(A\sqcup\{\ast\})^2$ has infinitely many elements outside it whenever $A$ is
infinite.  Thus the direct finite-subterm construction is needed if one wants
to remain inside the same class of completions.

We then use these completions to define finite-complement congruences on
$\Free D$.  For $m\ge 1$, write $x\fceq{m}y$ when every finite-complement
completion with at most $m$ new elements identifies $x$ and $y$.  The
finite-subset theorem implies
\[
 \bigcap_{m\ge1}\fceq{m}\;=\;\Delta_{\Free D}.
\]
Hence these congruences define a separated uniform structure, and we consider
its Cauchy completion.

The second part of the paper gives two criteria ensuring that this completion
contains new points.  For finite $A$, properness is equivalent to genuine
partiality of $D$.  For arbitrary $A$, we give a sufficient condition using a
unary context that fixes every base element.  The proof is again elementary:
finite-state periodicity produces a Cauchy subsequence, while a finite-subterm
separator rules out convergence to any element of $\Free D$.

The scope of the contribution is deliberately narrow.  Free completions of
partial algebras, residuality, discrimination, finite-state periodicity and
Cauchy completion are established ideas.  The point studied here is their
interaction under the additional requirement that a possibly infinite base be
retained pointwise while only its complement is finite.  The literature
comparison in \cref{sec:related-work} is written accordingly.

All proofs are given in ordinary mathematical language.  A machine-checked
formalization is mentioned only in an appendix as supplementary verification
and is not used as a premise in any argument.
